% Book 4, Year 2, chapter 17 -- The link layer.
% Provenance (never printed): sources/book4/BRIEFS.md, chapter 17.
\chapter{The Link Layer}\label{ch:b2:link-layer}

The town's lighthouse is told by radio when to light its lamp. One evening a \omterm{def:b2:link-layer:burst}{burst} of noise flips a single bit of
the message on its way:

\omoutput{code/bachelor-2/17-link-layer/outputs/lighthouse.txt}{One flipped bit, and who notices
it.}\label{out:b2:link-layer:hook}

``20:00'' became ``21:00'': the character \code{0} is the byte \code{0x30} and \code{1} is \code{0x31}, one bit
apart. The message still reads well, so no person would notice. The receiver noticed because the sender had added
four bytes computed from the message, a \emph{check}, and the received message no longer matches it. This chapter is
about the bottom of the network: how a stream of bits on one wire or one radio channel is cut into \omterm{def:b2:link-layer:frame}{frames}, how errors
are detected, how many machines share one medium, and how a \omterm{def:b2:link-layer:frame}{frame} finds its way across a local network.

\begin{recall}
Book 2: the network is built in layers, each using the one below; a protocol is a set of rules that both sides
follow. The physics series: a signal on a wire or in the air weakens with distance and picks up noise.
\cref{ch:b2:sockets}: bytes sent over a \omterm{def:b2:sockets:socket}{socket} arrive intact and in order; this chapter and the next two show who
makes that true.
\end{recall}

\section{Links and frames}\label{sec:b2:link-layer:frames}

\begin{definition}[Link, frame]\label{def:b2:link-layer:frame}
A \emph{link}\index{link} connects machines that can reach each other directly: a cable between two machines, or a
radio channel that all of them hear. The \emph{link layer}\index{link layer} carries \emph{frames}\index{frame}: a
frame is a block of bytes with a header (who sends it, to whom, what it carries), a payload, and a check. The link
layer hands frames to the layer above; a frame never leaves its link.
\end{definition}

The first problem is to find where a \omterm{def:b2:link-layer:frame}{frame} begins and ends in a stream of bits. A common answer marks the ends with a
special byte, the \emph{flag}. A flag byte inside the payload must then be disguised.

\begin{method}[Byte stuffing]\label{met:b2:link-layer:stuffing}
% ledger: b2:ppp-flag
\omterm{def:b2:link-layer:frame}{Frames} begin and end with the flag \code{0x7e}. Inside a \omterm{def:b2:link-layer:frame}{frame}, the sender replaces each byte \code{0x7e} or
\code{0x7d} by the escape \code{0x7d} followed by the byte XOR \code{0x20}. The receiver cuts the stream at the
flags, and in each \omterm{def:b2:link-layer:frame}{frame} replaces \code{0x7d}, $b$ by $b$ XOR \code{0x20}.
\end{method}

\omcode{code/bachelor-2/17-link-layer/listings/stuffing.py}{5}{27}{Stuffing and unstuffing}

\omoutput{code/bachelor-2/17-link-layer/outputs/stuffing.txt}{Two \omterm{def:b2:link-layer:frame}{frames} on the wire.}

After stuffing, \code{0x7e} occurs only at the ends of a \omterm{def:b2:link-layer:frame}{frame}, so the cut at the flags is exact. XOR with
\code{0x20} is its own inverse, which is why the receiver applies the same operation.

\section{Detecting errors}\label{sec:b2:link-layer:errors}

Noise flips bits.

\begin{definition}[Error-detecting code, checksum]\label{def:b2:link-layer:edc}
An \emph{error-detecting code}\index{error-detecting code} adds to each block of data a few check bits computed from
it; the receiver computes them again from what arrived, and drops the block if they differ. A \emph{checksum}
\index{checksum} is a check computed by adding the data as numbers; a \omterm{def:b2:link-layer:parity}{parity bit}, below, is the smallest check.
\end{definition}

No check catches every error: the question is which errors each one is sure to catch.

\subsection{Parity}

\begin{definition}[Parity bit]\label{def:b2:link-layer:parity}
The \emph{even parity bit}\index{parity bit} of a block of bits is 1 if the block has an odd number of 1s, and 0
otherwise, so that the block and its parity bit together hold an even number of 1s.
\end{definition}

\begin{proposition}[Parity]\label{prop:b2:link-layer:parity}
A \omterm{def:b2:link-layer:parity}{parity bit} detects every error that flips an odd number of bits of the block and its \omterm{def:b2:link-layer:parity}{parity bit}, and no error that
flips an even number.
\end{proposition}

\begin{proof}
Sent, the block and its bit hold an even number of 1s. Each flip changes that number by $\pm 1$, so after $k$ flips
it has changed by $k$ modulo 2: it is odd, and the receiver sees the error, exactly when $k$ is odd.
\end{proof}

\omoutput{code/bachelor-2/17-link-layer/outputs/parity.txt}{One \omterm{def:b2:link-layer:parity}{parity bit} per byte, then a grid of
parities.}\label{out:b2:link-layer:parity}

Arranged as a grid, with a \omterm{def:b2:link-layer:parity}{parity bit} per row (byte) and a parity byte for the columns, parity does more: a single
flip makes exactly one row and one column wrong, and their crossing is the bit to flip back. The code
\emph{corrects} the error (\cref{exo:b2:link-layer:3}). The price is one bit in nine, plus a byte per block.

\subsection{The Internet checksum}

% ledger: b2:inet-checksum
The Internet's own protocols use a cheaper check.

\begin{definition}[Internet checksum]\label{def:b2:link-layer:inet-checksum}
The \emph{Internet checksum}\index{Internet checksum} adds the data as 16-bit words, adding each carry out of the top
back in at the bottom (the \emph{one's complement sum}), and sends the complement of the result. The receiver adds
everything, \omterm{def:b2:link-layer:edc}{checksum} included, and expects \code{ffff}.
\end{definition}

\omcode{code/bachelor-2/17-link-layer/listings/inet_checksum.py}{4}{15}{The one's complement sum and the
\omterm{def:b2:link-layer:edc}{checksum}}

\omoutput{code/bachelor-2/17-link-layer/outputs/inet_checksum.txt}{It catches a flipped bit, not two swapped
words.}

A sum does not depend on the order of its terms, so swapping two words passes unseen. The \omterm{def:b2:link-layer:edc}{checksum} is fast to compute
in software, which is why the Internet layers use it; the \omterm{def:b2:link-layer:frame}{links}, which see the noise, use a stronger check.

\subsection{Cyclic redundancy codes}

A \omterm{def:b2:link-layer:crc}{cyclic redundancy code} treats a string of bits as a polynomial whose coefficients are 0 or 1: the bits $1101$ are
$x^3 + x^2 + 1$. Coefficients are added modulo 2, so adding is XOR, with no carries, and $x^k + x^k = 0$. Polynomials
can be multiplied and divided with remainder exactly as with numbers, the subtractions in a long division being XORs.

\begin{definition}[Cyclic redundancy code]\label{def:b2:link-layer:crc}
Let $g(x)$ be a polynomial of degree $r$ whose constant term is 1, the \emph{generator}. A message $M(x)$ is sent as
$T(x) = M(x)\,x^r + R(x)$, where $R(x)$ is the remainder of $M(x)\,x^r$ divided by $g(x)$: the message followed by
$r$ \emph{check bits}. The receiver accepts a \omterm{def:b2:link-layer:frame}{frame} $T'(x)$ if $g(x)$ divides it. This is a \emph{cyclic redundancy
code}\index{cyclic redundancy code}, or CRC.
\end{definition}

Since $M x^r = Q g + R$ and $R + R = 0$, $T = M x^r + R = Q g$: every \omterm{def:b2:link-layer:frame}{frame} sent is a multiple of $g$. The long
division is mechanical:

\omoutput{code/bachelor-2/17-link-layer/outputs/long_division.txt}{A CRC by long division, with $g(x) = x^4 + x +
1$: XOR $g$ under every leading 1.}

\omcode{code/bachelor-2/17-link-layer/listings/crc.py}{5}{12}{The same division in code, on the bits of an
integer}\label{lst:b2:link-layer:remainder}

If the noise adds an error $E(x)$ (a 1 at each flipped bit), the receiver gets $T + E$, and since $g$ divides $T$, it
accepts exactly when $g$ divides $E$. So the errors a CRC misses are the multiples of its generator, and choosing a
CRC is choosing $g$.

\begin{definition}[Burst error]\label{def:b2:link-layer:burst}
A \emph{burst}\index{burst error} of length $b$ is an error whose first and last flipped bits are $b$ positions apart,
counting both; the bits between them may be flipped or not.
\end{definition}

\begin{theorem}[Bursts]\label{thm:b2:link-layer:crc-burst}
A CRC whose generator has degree $r$ and constant term 1 detects every \omterm{def:b2:link-layer:burst}{burst} of length $b \le r$.
\end{theorem}

\begin{proof}
Such an error is $E(x) = x^i B(x)$, where $B$ has degree $b - 1 < r$ and constant term 1. First, if $g$ divides $x
C(x)$ for some $C$, then $g$ divides $C$: write $xC = gQ$; the constant term of $gQ$ is that of $Q$, because $g$'s is
1, and it is 0, because $xC$ has none; so $Q = xQ'$, and $xC = x\,gQ'$ gives $C = gQ'$. Applying this $i$ times, if
$g$ divided $E$ it would divide $B$. But a nonzero multiple of $g$ has degree at least $r$, and $B$ is nonzero of
degree less than $r$. So $g$ does not divide $E$.
\end{proof}

\begin{proposition}[Odd numbers of flips]\label{prop:b2:link-layer:crc-odd}
If $x + 1$ divides the generator, the CRC detects every error with an odd number of flipped bits.
\end{proposition}

The proof is \cref{exo:b2:link-layer:8}; the idea is to evaluate polynomials at $x = 1$. The 16-bit generator
% ledger: b2:crc16-ccitt
$x^{16} + x^{12} + x^5 + 1$ has four terms, so it is divisible by $x + 1$, and both results apply. They can be checked
by brute force on a \omterm{def:b2:link-layer:frame}{frame} of 4 bytes and its 16 check bits:

\omoutput{code/bachelor-2/17-link-layer/outputs/burst_check.txt}{Every \omterm{def:b2:link-layer:burst}{burst} of up to 17 bits through
CRC-16.}\label{out:b2:link-layer:bursts}

Not one \omterm{def:b2:link-layer:burst}{burst} of 16 bits or fewer passes. Of the \omterm{def:b2:link-layer:burst}{bursts} of 17 bits, exactly one pattern per position passes: the
generator itself, shifted.

\subsection{CRC-32 in practice}

% ledger: b2:crc32-ieee
The 32-bit CRC of Ethernet uses a generator of degree 32 written \code{0x104c11db7} as a bit vector. Real CRCs add two
details to the definition. The hardware sends each byte lowest bit first, so the bits are processed in that order,
and the generator appears reversed, as \code{0xedb88320}. The register starts at all ones, so that leading zero bytes
% ledger: b2:crc32-gzip
change the check, and the result is complemented. The same CRC-32 protects zip and gzip files and PNG images:

\omcode{code/bachelor-2/17-link-layer/listings/crc.py}{35}{58}{CRC-32 bit by bit, then a byte at a time from a
table}\label{lst:b2:link-layer:crc32}

\omoutput{code/bachelor-2/17-link-layer/outputs/crc.txt}{Three ways, one value: ours agrees with
\code{zlib.crc32}.}

The table holds, for each byte value, the effect of eight steps of the bitwise loop; one lookup replaces eight steps.
The same table in C gives the same values:

\omcode{code/bachelor-2/17-link-layer/listings/crc32.c}{18}{24}{The table-driven CRC-32 in C}

\omoutput{code/bachelor-2/17-link-layer/outputs/crc32.txt}{The values of the C program.}

\section{Sharing a medium}\label{sec:b2:link-layer:sharing}

On a radio channel, every station hears every other.

\begin{definition}[Multiple access protocol, collision]\label{def:b2:link-layer:mac-protocol}
When several stations share one medium and two transmit at once, their signals mix and both \omterm{def:b2:link-layer:frame}{frames} are lost: a
\emph{collision}\index{collision}. A \emph{multiple access protocol}\index{multiple access protocol} (or medium
access protocol) is the rule by which each station decides when to transmit.
\end{definition}

\begin{history}
% ledger: b2:aloha-1970
In 1970, Norman Abramson described the \omterm{def:b2:link-layer:aloha}{ALOHA} system, which linked terminals on several Hawaiian islands to one
computer over a shared radio channel. Its rule was the simplest possible: send when you have a \omterm{def:b2:link-layer:frame}{frame}; if no
acknowledgement comes, wait a random time and send again.
\end{history}

In \emph{slotted \omterm{def:b2:link-layer:aloha}{ALOHA}}, time is cut into slots of one \omterm{def:b2:link-layer:frame}{frame} each, and a station sends only at the start of a slot.
A slot is useful when exactly one station sends. Let $n$ stations each send in a slot with probability $p = G/n$, so
that $G$ \omterm{def:b2:link-layer:frame}{frames} are tried per slot on average.

\begin{proposition}[Slotted ALOHA]\label{prop:b2:link-layer:aloha}
The fraction of useful slots is $S_n = G(1 - G/n)^{n-1}$, which tends to $S = G e^{-G}$ as $n$ grows. The maximum of
$G e^{-G}$ is $1/e \approx 0.368$, at $G = 1$.
\end{proposition}

\begin{proof}
Exactly one station sends with probability $n \cdot p (1 - p)^{n-1} = G (1 - G/n)^{n-1}$: one of the $n$ sends, and
the $n - 1$ others do not. As $n$ grows, $(1 - G/n)^{n} \to e^{-G}$ and $(1 - G/n)^{-1} \to 1$. The maximum is
\cref{exo:b2:link-layer:10}.
\end{proof}

\omoutput{code/bachelor-2/17-link-layer/outputs/aloha_sim.txt}{Slotted \omterm{def:b2:link-layer:aloha}{ALOHA} with 100 stations, 20\,000 slots per
value of $G$.}\label{out:b2:link-layer:aloha}

\begin{omfigure}
\begin{tikzpicture}
\begin{axis}[omrecord, width=0.66\linewidth, height=5cm, xlabel={$G$, frames tried per slot},
    ylabel={useful slots}, xmin=0, xmax=3.1, ymin=0, ymax=0.42, legend pos=north east,
    legend style={font=\scriptsize}]
  \addplot[omDef, thick, domain=0:3.1, samples=80] {x*exp(-x)};
  \addlegendentry{$G e^{-G}$}
  \addplot[only marks, mark=*, mark size=1.6pt, omThm] table[x=G, y=simulated]
    {code/bachelor-2/17-link-layer/figures/aloha.dat};
  \addlegendentry{simulated}
  \draw[dashed, gray] (axis cs:1,0) -- (axis cs:1,0.368) -- (axis cs:0,0.368);
\end{axis}
\end{tikzpicture}
\par\smallskip{\small The points of \cref{out:b2:link-layer:aloha} against $G e^{-G}$. Too few attempts leave
slots empty; too many waste them in \omterm{def:b2:link-layer:mac-protocol}{collisions}. At best, 37 slots in 100 are useful.}
\end{omfigure}

On a cable, a station can do better: listen first. \emph{Carrier sense multiple access with \omterm{def:b2:link-layer:mac-protocol}{collision} detection}
(\omterm{def:b2:link-layer:aloha}{CSMA/CD}), the rule of the first Ethernets, sends only when the medium is quiet, keeps listening while sending, and
stops as soon as it hears a \omterm{def:b2:link-layer:mac-protocol}{collision}. It then waits a random number of slots before trying again, chosen by
\omterm{def:b2:link-layer:aloha}{binary exponential backoff} (defined below): after its $c$-th \omterm{def:b2:link-layer:mac-protocol}{collision} on a \omterm{def:b2:link-layer:frame}{frame}, a station
waits a number of slots drawn uniformly from $0$ to $2^{\min(c,10)} - 1$. Each \omterm{def:b2:link-layer:mac-protocol}{collision} doubles the range, so the
stations spread out until one gets through; \cref{exo:b2:link-layer:9} measures the delay.

\begin{definition}[ALOHA, CSMA/CD, binary exponential backoff]\label{def:b2:link-layer:aloha}
In \emph{ALOHA}\index{ALOHA} a station sends whenever it has a \omterm{def:b2:link-layer:frame}{frame} and, after a \omterm{def:b2:link-layer:mac-protocol}{collision}, again after a random
delay; \emph{slotted ALOHA} sends only at the start of a time slot. \emph{CSMA/CD}\index{CSMA/CD} senses the medium
before sending and stops on hearing a \omterm{def:b2:link-layer:mac-protocol}{collision}. \emph{Binary exponential backoff}\index{binary exponential backoff} draws the delay after the $c$-th
\omterm{def:b2:link-layer:mac-protocol}{collision} uniformly among $0, \dots, 2^{\min(c,10)} - 1$ slots.
\end{definition}

Carrier sense fails on radio, where a station may not hear every other one. If $A$ and $C$ are both in range of $B$ but
not of each other, each senses a quiet medium and sends to $B$, and their \omterm{def:b2:link-layer:frame}{frames} collide at $B$ without either hearing
it.

\begin{definition}[Hidden terminal]\label{def:b2:link-layer:hidden}
A \emph{hidden terminal}\index{hidden terminal} is a station whose transmissions reach a receiver but not another
sender to that receiver, so that carrier sense cannot prevent their \omterm{def:b2:link-layer:mac-protocol}{collisions}. Wireless networks therefore detect
lost \omterm{def:b2:link-layer:frame}{frames} by the missing acknowledgement, and may reserve the medium first by a short exchange.
\end{definition}

\section{Ethernet and switches}\label{sec:b2:link-layer:ethernet}

\begin{definition}[MAC address, broadcast address]\label{def:b2:link-layer:mac-address}
% ledger: b2:mac-bits
A \emph{MAC address}\index{MAC address} is the 6-byte address of a network interface on its \omterm{def:b2:link-layer:frame}{link}, written as six
hexadecimal bytes, \code{02:00:00:00:00:01}. In the first byte, the lowest bit marks a \emph{group} address (a \omterm{def:b2:link-layer:frame}{frame}
for several receivers), and the next bit a \emph{locally administered} one, chosen by its owner rather than by the
maker. The address \code{ff:ff:ff:ff:ff:ff}, the \emph{broadcast address}\index{broadcast address}, means every
interface on the \omterm{def:b2:link-layer:frame}{link}.
\end{definition}

% ledger: b2:ethernet-header
% ledger: b2:ethernet-payload
% ledger: b2:ethertypes
An Ethernet \omterm{def:b2:link-layer:frame}{frame} has a 14-byte header (destination, source, type), a payload of 46 to 1500 bytes, shorter payloads
being padded, and a 4-byte CRC-32. The type says what the payload is: \code{0800} an IPv4 packet, \code{0806} an ARP
message (\cref{sec:b2:link-layer:arp}), \code{86dd} an IPv6 packet.

\begin{omfigure}
\begin{tikzpicture}[font=\scriptsize, f/.style={draw, minimum height=7mm, inner sep=2pt, align=center}]
\node[f, fill=omDef!25, minimum width=2.1cm] (d) at (0,0) {destination\\ 6 bytes};
\node[f, fill=omDef!25, minimum width=2.1cm, anchor=west] (s) at (d.east) {source\\ 6 bytes};
\node[f, fill=omDef!25, minimum width=1.2cm, anchor=west] (t) at (s.east) {type\\ 2};
\node[f, fill=omThm!10, minimum width=4.6cm, anchor=west] (p) at (t.east) {payload\\ 46 to 1500 bytes};
\node[f, fill=omDef!10, minimum width=1.5cm, anchor=west] (c) at (p.east) {CRC-32\\ 4 bytes};
\draw[decorate, decoration={brace, mirror, amplitude=4pt}] ([yshift=-1mm]d.south west) -- node[below=5pt]
{header, 14 bytes} ([yshift=-1mm]t.south east);
\draw[decorate, decoration={brace, mirror, amplitude=4pt}] ([yshift=-1mm]p.south west) -- node[below=5pt]
{a packet of the layer above, padded to 46} ([yshift=-1mm]p.south east);
\end{tikzpicture}
\par\smallskip{\small An Ethernet \omterm{def:b2:link-layer:frame}{frame}. The CRC-32 covers the header and the payload; a \omterm{def:b2:link-layer:frame}{frame} is therefore 64 to
1518 bytes long.}
\end{omfigure}

\omcode{code/bachelor-2/17-link-layer/listings/ethernet_frame.py}{13}{26}{A \omterm{def:b2:link-layer:frame}{frame} built byte by
byte}\label{lst:b2:link-layer:frame}

\omoutput{code/bachelor-2/17-link-layer/outputs/ethernet_frame.txt}{The 64 bytes of a broadcast \omterm{def:b2:link-layer:frame}{frame} carrying an
ARP request.}\label{out:b2:link-layer:frame}

\begin{definition}[Network switch, switch table]\label{def:b2:link-layer:switch}
A \emph{network switch}\index{network switch} is a box with many ports, each linked to a machine or another switch,
that receives each \omterm{def:b2:link-layer:frame}{frame} and sends it on only through the ports where its destination may be. Its \emph{switch table}
maps MAC addresses to ports.
\end{definition}

Machines on one Ethernet are connected by switches. A switch needs no configuration, because it learns its table.

\begin{method}[Self-learning switch]\label{met:b2:link-layer:switch}
For each \omterm{def:b2:link-layer:frame}{frame} that arrives on port $p$: record in the table that its source is reachable through $p$. Then, if the
destination is in the table on a port other than $p$, send the \omterm{def:b2:link-layer:frame}{frame} on that port only (\emph{forward}); if it is
in the table on port $p$, drop it (\emph{filter}: the receiver already heard it); if it is not in the table, or is a
broadcast, send it on every port except $p$ (\emph{flood}).
\end{method}

\omcode{code/bachelor-2/17-link-layer/listings/switch_sim.py}{4}{16}{A self-learning switch}

\omoutput{code/bachelor-2/17-link-layer/outputs/switch_sim.txt}{Five hosts on ports 1 to 5; the table after each
\omterm{def:b2:link-layer:frame}{frame}.}

The first \omterm{def:b2:link-layer:frame}{frame} from $A$ is flooded, since nobody knows $B$ yet; $B$'s answer is forwarded, since $A$ is now known.
After a few \omterm{def:b2:link-layer:frame}{frames} the switch floods only for unknown or broadcast destinations. Several switches can be connected in
a tree, and learning works through them (\cref{exo:b2:link-layer:5}). A loop, though, would make a flooded \omterm{def:b2:link-layer:frame}{frame}
circle forever, each switch flooding it again. Real networks keep a spare cable for safety anyway; the switches then
run the \emph{spanning tree protocol}, in which they elect one switch as root and switch off the ports that would
close a loop, leaving a tree, and switch a spare port on again when a cable fails.

\omimg{book4/ethernet-rj45-connector.jpg}{0.42\linewidth}{The plug of an Ethernet cable: eight gold contacts for
four twisted pairs of copper wires. Photograph: David Monniaux, CC BY-SA 3.0.}

\section{Address resolution}\label{sec:b2:link-layer:arp}

The layer above addresses machines by IP address (\cref{ch:b2:network-layer}); the \omterm{def:b2:link-layer:frame}{link} needs a \omterm{def:b2:link-layer:mac-address}{MAC address}.

\begin{definition}[Address resolution]\label{def:b2:link-layer:arp}
\emph{Address resolution}\index{address resolution protocol} finds the \omterm{def:b2:link-layer:mac-address}{MAC address} of the interface that holds a given
IP address on the same \omterm{def:b2:link-layer:frame}{link}; on Ethernet with IPv4 it is done by the \emph{address resolution protocol} (ARP).
\end{definition}

\begin{method}[ARP]\label{met:b2:link-layer:arp}
% ledger: b2:arp-fields
To send to IP address $a$ on its \omterm{def:b2:link-layer:frame}{link}, a machine looks $a$ up in its \emph{ARP cache}. If the address is missing,
it broadcasts a request (operation 1): ``who has $a$? tell me, at this IP and \omterm{def:b2:link-layer:mac-address}{MAC address}''. The machine whose
address is $a$ records the asker in its cache and sends a reply (operation 2) with its \omterm{def:b2:link-layer:mac-address}{MAC address}; the asker records
it. Entries expire after a while, so that a machine that changed its interface is found again.
\end{method}

The request of \cref{out:b2:link-layer:frame} is such a broadcast: hardware type 1 (Ethernet), protocol
\code{0800} (IPv4), address lengths 6 and 4, operation 1, the sender's two addresses, an unknown target \omterm{def:b2:link-layer:mac-address}{MAC address}
(zeros) and the target IP address \code{0a 00 00 02}, that is 10.0.0.2.

\omoutput{code/bachelor-2/17-link-layer/outputs/arp_sim.txt}{Three hosts; cache entries live 60 seconds in this
simulation.}

\begin{omfigure}
\begin{tikzpicture}[font=\scriptsize, w/.style={-{Stealth[length=1.8mm]}, thick},
  l/.style={fill=white, inner sep=1.5pt}]
\foreach \x/\n in {0/A, 4.5/B, 9/C} {\node[font=\small\bfseries] at (\x,0.3) {\n}; \draw[gray] (\x,0) --
(\x,-3.6);}
\draw[w, omThm] (0,-0.6) -- (9,-0.6);
\node[l, text=omThm, above=1.5pt] at (4.5,-0.6) {broadcast: who has 10.0.0.3? tell 10.0.0.1};
\node[anchor=west] at (4.6,-1.05) {not mine: ignored};
\node[anchor=east] at (8.9,-1.5) {records A in its cache};
\draw[w, omDef] (9,-2.1) -- (0,-2.1);
\node[l, text=omDef, above=1.5pt] at (4.5,-2.1) {reply to A only: 10.0.0.3 is at 02:00:00:00:00:03};
\node[anchor=west] at (0.1,-2.55) {records C in its cache};
\draw[w, black] (0,-3.3) -- (9,-3.3);
\node[l, above=1.5pt] at (4.5,-3.3) {the frame, sent to C's MAC address};
\end{tikzpicture}
\par\smallskip{\small An ARP exchange. The request reaches everyone; the reply and the \omterm{def:b2:link-layer:frame}{frame} reach only their
addressee. A second \omterm{def:b2:link-layer:frame}{frame} to C needs no request until the entry expires.}
\end{omfigure}

ARP trusts every reply. A machine that answers for an address that is not its own receives the \omterm{def:b2:link-layer:frame}{frames} meant for it;
networks defend against this by having switches check replies against the addresses they handed out, and by
encrypting what travels above the \omterm{def:b2:link-layer:frame}{link}, as the Year 3 volume explains.

\section{Wireless links}\label{sec:b2:link-layer:wireless}

A radio \omterm{def:b2:link-layer:frame}{link} differs from a cable in three ways. The signal weakens quickly with distance and through walls, so its
error rate is far higher, and many wireless \omterm{def:b2:link-layer:frame}{links} resend a \omterm{def:b2:link-layer:frame}{frame} themselves: the receiver acknowledges each \omterm{def:b2:link-layer:frame}{frame},
and the sender repeats it when no acknowledgement comes. Second, a station cannot listen while it sends, its own
signal drowning every other, so it cannot detect a \omterm{def:b2:link-layer:mac-protocol}{collision}. Third, not every station hears every other: these are the
\omterm{def:b2:link-layer:hidden}{hidden terminals} of \cref{def:b2:link-layer:hidden}, whose \omterm{def:b2:link-layer:frame}{frames} collide at the receiver. Wireless networks
therefore \emph{avoid} \omterm{def:b2:link-layer:mac-protocol}{collisions} (CSMA/CA): a station waits a random backoff even when the channel is quiet, and
may first reserve the channel with a short request that the receiver answers with a short ``clear to send'', which
every station near the receiver, hidden or not, hears.

\omimg{book4/ai/b2-link-layer-plugging-cable.jpg}{0.58\linewidth}{Plugging a cable into a switch. Each port is a \omterm{def:b2:link-layer:frame}{link}
of its own: the switch decides which \omterm{def:b2:link-layer:frame}{links} a \omterm{def:b2:link-layer:frame}{frame} crosses.}

\section{Exercises}\label{sec:b2:link-layer:exercises}

\begin{exercise}[$\star$]\label{exo:b2:link-layer:1}
Give the \omterm{def:b2:link-layer:parity}{even parity bits} of the bytes \code{0x4c}, \code{0x41}, \code{0x4d} and \code{0x50}, and check your answer
against \cref{out:b2:link-layer:parity}. Which of the following errors does a \omterm{def:b2:link-layer:parity}{parity bit} per byte detect: one flip in
each of two bytes; three flips in one byte; two flips in one byte?
\end{exercise}

\begin{exercise}[$\star$]\label{exo:b2:link-layer:2}
This \omterm{def:b2:link-layer:edc}{checksum} forgets to add its carries back. Find the line, and explain why the first input gives the right value
and the second does not.
\omcode{code/bachelor-2/17-link-layer/listings/exo_nofold.py}{5}{9}{No folding}
\omoutput{code/bachelor-2/17-link-layer/outputs/exo_nofold.txt}{Right and wrong.}
\end{exercise}

\begin{exercise}[$\star$]\label{exo:b2:link-layer:3}
Write \code{correct(data, rows, col)}, which takes the data, its row \omterm{def:b2:link-layer:parity}{parity bits} and its column parity byte, corrects
any single flipped bit (in the data or in the parities) and returns \code{None} when it cannot.
\end{exercise}

\begin{exercise}[$\star$]\label{exo:b2:link-layer:4}
By long division, compute the 3 check bits of the message $1101\,0011$ with $g(x) = x^3 + x + 1$. Check that the \omterm{def:b2:link-layer:frame}{frame}
sent is divisible by $g$.
\end{exercise}

\begin{exercise}[$\star\star$]\label{exo:b2:link-layer:5}
Two switches are connected by a cable between port 3 of $s_1$ and port 1 of $s_2$; hosts $A$ and $B$ are on ports 1
and 2 of $s_1$, $C$ and $D$ on ports 2 and 3 of $s_2$. Starting from empty tables, the \omterm{def:b2:link-layer:frame}{frames} $A \to C$, $C \to A$,
$D \to B$, $B \to D$, $A \to B$, $D \to C$ are sent. Give both tables at the end, and say which \omterm{def:b2:link-layer:frame}{frames} were flooded
by which switch.
\end{exercise}

\begin{exercise}[$\star\star$]\label{exo:b2:link-layer:6}
This decoder drops the escape bytes. Explain its output and fix it.
\omcode{code/bachelor-2/17-link-layer/listings/exo_unstuff.py}{5}{6}{A wrong decoder}
\omoutput{code/bachelor-2/17-link-layer/outputs/exo_unstuff.txt}{What came out.}
\end{exercise}

\begin{exercise}[$\star\star$]\label{exo:b2:link-layer:7}
Write a CRC with the generator $x^8 + x^2 + x + 1$, most significant bit first, bit by bit and with a table of 256
entries, and test that they agree with each other and with \code{remainder} (\cref{lst:b2:link-layer:remainder}).
\end{exercise}

\begin{exercise}[$\star\star$]\label{exo:b2:link-layer:8}
Prove \cref{prop:b2:link-layer:crc-odd}. Hint: for a polynomial $P$, $P(1)$ is the number of its terms modulo 2.
\end{exercise}

\begin{exercise}[$\star\star$]\label{exo:b2:link-layer:9}
Simulate $k$ stations that each have one \omterm{def:b2:link-layer:frame}{frame} to send at time 0 and follow \omterm{def:b2:link-layer:aloha}{binary exponential backoff} on a slotted
channel. Measure the mean number of slots until every \omterm{def:b2:link-layer:frame}{frame} is through and the mean delay of a \omterm{def:b2:link-layer:frame}{frame}, for $k = 2, 4,
\dots, 32$. How do they grow with $k$?
\end{exercise}

\begin{exercise}[$\star\star\star$]\label{exo:b2:link-layer:10}
Prove that $G e^{-G}$ is largest at $G = 1$. In \emph{pure} \omterm{def:b2:link-layer:aloha}{ALOHA}, a station sends whenever it likes; a \omterm{def:b2:link-layer:frame}{frame} is lost
if any other \omterm{def:b2:link-layer:frame}{frame} starts within one \omterm{def:b2:link-layer:frame}{frame} time before or after it. With attempts forming a Poisson process of $G$
\omterm{def:b2:link-layer:frame}{frames} per \omterm{def:b2:link-layer:frame}{frame} time, the throughput is $G e^{-2G}$. Find its maximum, and explain the factor 2.
\end{exercise}

\begin{exercise}[$\star\star\star$]\label{exo:b2:link-layer:11}
Show that a \omterm{def:b2:link-layer:burst}{burst} $x^i B(x)$ passes a CRC exactly when $B$ does, so that checking every \omterm{def:b2:link-layer:burst}{burst} on a \omterm{def:b2:link-layer:frame}{frame} reduces to
checking every pattern $B$ once. For CRC-16, count the patterns of lengths 17 and 18 that pass, compare with
\cref{out:b2:link-layer:bursts}, and explain both counts.
\end{exercise}

\begin{exercise}[$\star\star\star$]\label{exo:b2:link-layer:12}
On \omterm{def:b2:link-layer:frame}{frames} of $N \ge r$ bits, show that exactly $2^{N-r} - 1$ of the $2^N - 1$ nonzero error patterns pass a CRC of
degree $r$, about one in $2^r$. With $x + 1$ dividing $g$, what fraction of the patterns with an even number of
flips pass?
\end{exercise}

\section{Problem: The noisy link to the lighthouse}\label{sec:b2:link-layer:problem}

\begin{problem}[{Weekend problem --- The noisy link to the lighthouse}]\label{pb:b2:link-layer:1}
The harbour sends short \omterm{def:b2:link-layer:frame}{frames} to the lighthouse over a radio \omterm{def:b2:link-layer:frame}{link}: a 4-byte command followed by 16 check bits of
CRC-16, with $g(x) = x^{16} + x^{12} + x^5 + 1$. Noise comes in \omterm{def:b2:link-layer:burst}{bursts}.

\needspace{4\baselineskip}\medskip\noindent\textbf{Part I --- Framing.}\par\nopagebreak
\begin{enumerate}
\item \omterm{def:b2:link-layer:frame}{Frames} are separated by the flag \code{0x7e}. Stuff the payload \code{7d 7e 00 7e}.
\item What is the longest a stuffed payload of $n$ bytes can be, and for which payloads?
\item Both flags between two \omterm{def:b2:link-layer:frame}{frames} are lost in the noise. What does the receiver see, and does the CRC of this
problem catch it?
\item Why does the CRC cover the payload before stuffing rather than after?
\item The lighthouse could instead receive a length byte before each payload. Why is a length byte risky on a noisy
\omterm{def:b2:link-layer:frame}{link}, where a flag is not?
\item How many bytes cross the \omterm{def:b2:link-layer:frame}{link} for the command \code{LAMP}, its check and two flags?
\end{enumerate}

\needspace{4\baselineskip}\medskip\noindent\textbf{Part II --- Detecting errors.}\par\nopagebreak
\begin{enumerate}[resume, beginpenalty=10000]
\item The last line of the output of \code{crc.py} gives the check bits of \code{LAMP}. Write the 6 bytes sent, and
say why the receiver's CRC of the whole \omterm{def:b2:link-layer:frame}{frame} is 0.
\item Why can one \omterm{def:b2:link-layer:parity}{parity bit} per \omterm{def:b2:link-layer:frame}{frame} not protect this \omterm{def:b2:link-layer:frame}{link} against \omterm{def:b2:link-layer:burst}{bursts}?
\item Using \cref{thm:b2:link-layer:crc-burst}, which \omterm{def:b2:link-layer:burst}{bursts} is CRC-16 sure to detect?
\item Using \cref{prop:b2:link-layer:crc-odd}, which other errors is it sure to detect?
\item In \cref{out:b2:link-layer:bursts}, why are there 32 undetected \omterm{def:b2:link-layer:burst}{bursts} of length 17 and not more?
\item The listing below draws $2^{24}$ random error patterns on the 48-bit \omterm{def:b2:link-layer:frame}{frame}. What count of undetected patterns
does \cref{exo:b2:link-layer:12} predict? Compare with the output.
\omoutput{code/bachelor-2/17-link-layer/outputs/random_errors.txt}{Random error patterns through CRC-16.}
\end{enumerate}

\needspace{4\baselineskip}\medskip\noindent\textbf{Part III --- Sharing the radio channel.}\par\nopagebreak
\begin{enumerate}[resume, beginpenalty=10000]
\item Five buoys report to the harbour on the same channel with slotted \omterm{def:b2:link-layer:aloha}{ALOHA}. Each sends in a slot with
probability $p$. What is the probability that a slot is useful?
\item Which $p$ makes it largest, and what is the largest value? Compare with $1/e$.
\item Why can the buoys not use \omterm{def:b2:link-layer:mac-protocol}{collision} detection?
\item The lighthouse and a buoy behind the headland both hear the harbour, but not each other. Name the problem and
say what it causes.
\item How does a short reservation exchange help, and what does it cost?
\item Why does every \omterm{def:b2:link-layer:frame}{frame} need an acknowledgement on this \omterm{def:b2:link-layer:frame}{link}, when a cable often needs none?
\end{enumerate}

\needspace{4\baselineskip}\medskip\noindent\textbf{Part IV --- Choosing a CRC.}\par\nopagebreak
\begin{enumerate}[resume, beginpenalty=10000]
\item \omterm{def:b2:link-layer:burst}{Bursts} of up to 20 bits are common on this \omterm{def:b2:link-layer:frame}{link}. Is CRC-16 enough? What does CRC-32 guarantee instead?
\item What does CRC-32 cost per \omterm{def:b2:link-layer:frame}{frame}, in bytes and in computation, against CRC-16?
\item A \omterm{def:b2:link-layer:frame}{frame} carries 4 bytes of data. What fraction of the \omterm{def:b2:link-layer:frame}{frame} is check bits with each CRC?
\item \omterm{def:b2:link-layer:burst}{Bursts} of 20 bits occur in one \omterm{def:b2:link-layer:frame}{frame} in a thousand. Estimate how many corrupted \omterm{def:b2:link-layer:frame}{frames} CRC-16 lets through per
million \omterm{def:b2:link-layer:frame}{frames}.
\item Why is a CRC computed by a table in software, or by a shift register in hardware, rather than by the
bit-by-bit division of \cref{lst:b2:link-layer:remainder}?
\item State the result: no \omterm{def:b2:link-layer:burst}{burst} of 16 bits or fewer passes CRC-16, and random errors pass at the predicted rate
$2^{-16}$, with the counts of the simulation.
\end{enumerate}
\end{problem}
